% Generated draft: requires mathematical and source review.
\documentclass[11pt]{article}
\usepackage[margin=25mm]{geometry}
\usepackage{fontspec}
\setmainfont{Latin Modern Roman}
\usepackage{amsmath,amssymb,mathtools,unicode-math}
\setmathfont{Latin Modern Math}
\usepackage{enumitem,needspace,xurl,hyperref}
\setlength{\parindent}{0pt}
\setlength{\parskip}{5pt}
\setlength{\emergencystretch}{3em}
\newcommand{\N}{\mathbb{N}}
\newcommand{\Z}{\mathbb{Z}}
\newcommand{\Q}{\mathbb{Q}}
\newcommand{\R}{\mathbb{R}}
\newcommand{\C}{\mathbb{C}}
\newcommand{\F}{\mathbb{F}}
\newcommand{\A}{\mathbb{A}}
\newcommand{\PP}{\mathbb{P}}
\newcommand{\E}{\mathbb{E}}
\newcommand{\eps}{\varepsilon}
\newcommand{\dd}{\,\mathrm{d}}
\begin{document}
\section*{UnsolvedMath ID: 30004601}
\textbf{Source problem code:} catalogue entry
\subsection{Definitions and mathematical statement}
A central hyperplane arrangement \(\mathcal A\) in \(K^\ell\) is a finite collection of hyperplanes through the origin, where \(K\) is a field. Put \(S=K[x_1,\ldots,x_\ell]\), with total-degree grading. If \(\alpha_H\) is a nonzero linear form defining \(H\in\mathcal A\), its defining polynomial is
\[
Q(\mathcal A)=\prod_{H\in\mathcal A}\alpha_H.
\]
The supplied record uses \(J(\mathcal A)\) without defining it. Two conventions requiring clarification are the ideal generated by the partial derivatives of \(Q(\mathcal A)\), and the ideal generated by \(Q(\mathcal A)\) together with those derivatives:
\[
J_{\partial}=\left(\frac{\partial Q}{\partial x_1},\ldots,\frac{\partial Q}{\partial x_\ell}\right),\qquad
J_{Q,\partial}=(Q)+J_{\partial}.
\]
These conventions agree when \(\deg Q\) is invertible in \(K\), by Euler's identity; the record supplies no characteristic hypothesis ensuring this.

The notation \(\operatorname{rgin}\) conventionally denotes the generic initial ideal for graded reverse lexicographic order, ordinarily with \(x_1>\cdots>x_\ell\): it is the monomial initial ideal obtained after a sufficiently general invertible linear change of coordinates. The record does not explicitly specify this order, a field hypothesis supporting the generic-coordinate construction over \(K\), or an extension-field convention. These choices must be confirmed rather than treated as additional source hypotheses.

With the intended meanings of \(J(\mathcal A)\) and \(\operatorname{rgin}\) fixed, write
\[
I=\operatorname{rgin}(J(\mathcal A)),\qquad
D=\{d\in\mathbb Z_{\geq 0}:x_1^d\in I\}.
\]
A minimal monomial generator of \(I\) is a monomial in its unique minimal monomial generating set, denoted \(G(I)\). A monomial \(t=x_1^{a_1}\cdots x_\ell^{a_\ell}\) involves \(x_3\) precisely when \(a_3>0\), equivalently \(x_3\mid t\); its degree is \(\deg t=\sum_i a_i\).

\textbf{Source conjecture, subject to the unresolved conventions above.} For every central arrangement in the intended class, set \(d_0=\min D\). Must
\[
\forall t\in G(I),\qquad x_3\mid t\ \Longrightarrow\ \deg t\geq d_0
\]
hold?

This formulation requires \(\ell\geq 3\) for \(x_3\) to exist and \(D\neq\varnothing\) for \(d_0\) to be defined. Neither requirement is explicitly stated in the supplied record. No value is assigned to \(\min\varnothing\), and no claim is made here that the missing conditions follow from the supplied hypotheses. The source does not specify whether empty arrangements or other degenerate cases are excluded.
\subsection{Short English statement}
The conjecture asks whether every minimal monomial generator involving the third variable has degree at least the least exponent of a pure power of the first variable in the generic initial ideal of the arrangement's Jacobian ideal. The field hypotheses, ideal and generic-initial-ideal conventions, and existence of that least exponent require clarification. The supplied registry labels the problem open; this draft does not independently verify that status.
\subsection{Sources}
\begin{itemize}
\item Supplied local UnsolvedMath registry record, permanent ID \(30004601\), source code \texttt{catalogue entry}, titled \emph{Degree Bounds for Generic Initial Ideals of Arrangements}. Registry locations identified by the project instructions are \texttt{data/unsolvedmath/problems.json} and \texttt{catalog/problems.json}; these files were not inspected for this draft. The record gives proposed year 2021 and names no individual proposer.
\item Supplied primary report citation: \emph{Logarithmic Vector Fields and Freeness of Divisors and Arrangements: New perspectives and applications} (2022), Oberwolfach Reports, pp.\ 225--300, DOI \texttt{10.4171/owr/2021/5}. \url{https://doi.org/10.4171/owr/2021/5}. The background describes these reports as documenting mathematical workshops and their open-problem discussions.
\item Supplied literature source: Bigatti--Palezzato--Torielli, \emph{Journal of Algebraic Combinatorics}. \url{https://eprints.lib.hokudai.ac.jp/repo/huscap/all/80515/J.%20algebr.%20comb%2051-2_297-315.pdf}.
\item Supplied literature source: Palezzato--Torielli, \emph{Journal of Algebra}. \url{https://eprints.lib.hokudai.ac.jp/dspace/bitstream/2115/91057/1/J.%20Algebra%20590_215-233.pdf}.
\item Supplied literature source: Oberwolfach Report 5/2021. \url{https://doi.org/10.4171/owr/2021/5}.
\item Literature assessment checked at \texttt{2026-08-22}: the supplied curation labels the problem open and reports that the generic-initial-ideal statement remains a conjecture in later work on Lefschetz properties, with no resolving proof or counterexample found. This is supplied dated curation derived from the OWR material, not an independently verified literature assessment; time-sensitive status and cited claims require independent verification.
\item The registry labels the statement \texttt{corrected\_verified} and reports manual semantic, typographic, and LaTeX review of its clean statement. Those supplied review labels do not establish mathematical correctness or resolve the formulation issues recorded here.
\end{itemize}
\end{document}
